% Written weak-base construction, pinned 2026-11-14 source snapshot.
% This fragment is designed to be included after term_leaf_completion.tex.
\section{A direct weak-base compiler from ordinary PA equality proofs}
\label{sec:forward-calibration-completion}

\paragraph{Statement and scope.}
Let $\mathcal C_4$ be the numerical calculus with the six unchanged
heads, the separate diagonal head $\Gamma(B)=I(B,U(B,B))$, the retained
guarded old gate, the dominated $\rho _8$ head, modus ponens, and theorem
generalization.  Use the exact numerical constructors and guard of the
machine, specified below.  There is a primitive recursive function $F$
such that
\begin{equation}
 \mathsf I\Sigma _1\ \vdash\quad
 \operatorname{OrdinaryPAeqProof}(d,\bot)\ \longrightarrow\
 \operatorname{C4Proof}(F(d),1).
 \tag{F1}
\end{equation}
Here the predicate on the left is the pinned upstream
\file{Presentation.ordinaryPAeq} proof predicate.  It is not an
unidentified substitute for the current Foundation PA predicate.  The
separate Foundation adapter must identify that endpoint before (F1) is
composed with the machine theorem.  The companion Foundation completion
provides that adapter; the present section supplies the ordinary-PAeq
to numerical-calculus part.

The proof below constructs direct derivation lists.  It neither invokes
the full-PA correctness theorem for the upstream \file{verifyValue}
checker nor asserts that that theorem has been weakened in Lean.
Source references specify finite constructor recipes and their displayed
formula equations.  We replace the upstream existential supplier
interfaces by the explicit finite-demand algorithms below.  This is a
written arithmetized proof, not a newly kernel-checked development.

All source paths in this section are relative to
\file{Computability/BusyBeaver/Research/PA/jamesyc-upstream/pa196/}
at commit
\file{1c9524d8bf2e0a27b9a16ff622359927f536e663}.
The accompanying source manifest records the exact blob identifiers.
Except where a longer prefix is shown, a filename in this section is
in \file{PA196/Arithmetic/Transport/}.

\subsection{The target and its primitive proof constructors}

Put
\[
 \pi(a,b)=\frac{(a+b)(a+b+1)}2+b,\qquad
 I(a,b)=4\pi(a,b),\quad M(a,b)=I(a,b)+1,
\]
\[
 U(a,b)=I(b,a)+2,\qquad E(a,b)=I(a,b)+3,
 \quad N(A)=I(A,1).
\]
Numerical conjunction is $N(I(A,N(B)))$ and
$\exists_a A=N(U(a,N(A)))$.  Code $0$ is truth and code $1$ is false;
in particular $I(0,0)=0$ is retained literally.  The six heads are
\begin{align*}
 h_1(D,H)&=I(\exists_D U(D,H),H),\\
 h_2(D,A,H)&=I(U(D,U(A,H)),U(A,U(D,H))),\\
 h_3(D,A,H)&=I(M(H,D),I(E(H,A),M(A,D))),\\
 h_4(D,A,H)&=I(E(H,D),I(M(A,H),M(A,D))),\\
 h_5(D,H)&=\exists_D E(H,D),\\
 h_6(D,A,H)&=I(U(D,A),I(U(D,I(A,H)),U(D,H))).
\end{align*}
The propositional head is
\[
 \rho _8(D,A,H,u)=I(I(I(H,u),D),I(I(D,H),I(A,H))),
 \qquad u\le D.
\]
For the old gate put
\[
 \begin{split}
 V&=I(M(a,b),U(c,M(a,b))),\qquad P=I(c,M(c,b)),\\
 S&=U(b,(I(M(a,b),P))\mathbin\wedge(I(M(c,b),c))),\\
 Q&=U(c,I(S,N(P))),\qquad G(a,b,c)=N(I(V,Q)).
 \end{split}
\]
It is a head when $c\ne a,b$; the machine's failed guard supplies
$0$.  The separate $\Gamma(B)=I(B,U(B,B))$ is always available.

A direct proof list records each numerical formula together with one
of: its head name and parameters, two earlier indices for MP, or one
earlier index and a label for Gen.  Checking a row uses only numerical
equalities, inequalities, and bounded earlier-index lookups.  The check
of the whole list is therefore primitive recursive.  Concatenation
with index relocation, appending MP or Gen, and a one-row head proof
are primitive recursive operations.  Their correctness is a bounded
Boolean induction on the list length.  These operations interpret the
upstream \file{Certificate} and \file{Context} recipes; no
\file{verifyValue} certificate code occurs in this interpretation.

The restricted propositional basis is sufficient.  Interpret
\file{M6/Logic/RestrictedImplication.lean}'s
\file{Rules}, \file{rho8All}, \file{composition},
\file{exchange}, \file{k}, \file{peirce}, and \file{s}
with the preceding list operations.  Every supplied restricted row
has its actual side condition $u\le D$: the only bound constructors
used by that recipe are reflexivity and
\[
 u\le x\Longrightarrow u\le I(x,y),\qquad
 u\le y\Longrightarrow u\le I(x,y).
\]
They follow from $x,y\le\pi(x,y)\le4\pi(x,y)$, including the
zero case.  These are the exact instances in
\file{PA196/Logic/ImplicationBasis.lean} and
\file{PA196/Arithmetic/Certificates/ImplicationBasis.lean}.
Thus identity, deduction, composition, exchange, contraposition,
double-negation elimination, and the finite classical connective
templates are uniform primitive recursive proof constructors.
Deduction itself traverses a supplied finite assumption/closed/MP
tree, using identity, $K$, and $S$.  Its correctness is induction on
that finite tree, not an assumption of a semantic deduction theorem.

\subsection{Replacing primitive equality vacuity at the exact boundary}

This replacement is necessary: $\mathcal C_4$ does not contain the old
primitive equality-vacuity axiom.  The old gate directly gives
membership vacuity and pointed separation by the fixed negated-
implication projections.  Generalized implication, negation, and
universal vacuity then use only those inputs and $h_1,h_2,h_6$.

Here are the finite replacement recipes, as written in the pinned
\file{PA/docs/pa194q_conpa_2026-11-13/references/calibration/}
files \file{diagonal_calculus.tex},
\file{separate_head_calculus.tex}, and
\file{original_pa195_calibration.tex}; these three paths are
relative to the repository root rather than the upstream prefix.
For $a\ne b$ let $\chi(t)$ be $M(d,t)$ or $M(t,d)$, with
$d\ne a$ in both of these fixed-label cases, or let it be $M(t,t)$.
Form
\[
 C=I(E(a,b),I(\chi(b),\chi(a))),\qquad \varphi=N(C),
 \qquad c=\varphi.
 \tag{F2}
\]
The label $c$ is the code of the complete negated conclusion, not a
label chosen by an unbounded freshness search.  Constructor
inequalities give $c>a,b$ and, in the first two cases, $c>d$.
Consequently the old gate at $(a,a,c)$ is licensed.  Its separation
matrix implies
\[
 S\longrightarrow U(a,I(M(c,a),N(\chi(a)))).
\]
The retained \file{negativeSpecialization} recipe uses only the
unchanged witness head, direct membership rewrites, and membership
vacuity, to obtain
\[
 S\longrightarrow I(M(c,b),N(\chi(b))).
\]
Under $\varphi$, the point condition gives $M(c,a)$, direct rewriting
gives $M(c,b)$, and this contradicts $\chi(b)$.  Hence
$\varphi\to(S\to\neg P)$.  Gen at $c$, distribution, and the
literal diagonal head $\Gamma(\varphi)$ give $\varphi\to Q$;
the same gate gives $\neg Q$.  Double-negation elimination proves
$C$.  There is no step generalizing an equality-containing assumption
at an arbitrary fresh label.  The cases $a=b$ are propositional; the
cases $d=a$ use the diagonal $\chi(t)=M(t,t)$ case and the two direct
rewrites.  This supplies both reverse membership directions uniformly
for every coincidence pattern of labels.

Define the neighborhood equality
\[
 R(a,b,z)=U(z,M(a,z)\leftrightarrow M(b,z))
          \wedge U(z,M(z,a)\leftrightarrow M(z,b)).
\]
Set $z_0=E(a,b)>a,b$.  The head $\Gamma(E(a,b))$ and the preceding
two-direction rewrites give $E(a,b)\to R(a,b,z_0)$.  For
$z\ne a,b,z_0$, change each neighborhood binder from $z_0$ to $z$
using the unchanged equality witness, both membership rewrites, and
membership-only vacuity.  If $z=z_0$, use identity.  This yields
\[
 E(a,b)\to R(a,b,z)\quad(z\ne a,b),\qquad
 \exists_a R(a,b,z)\quad(z\ne a,b).
 \tag{F3}
\]
All operations in (F2)--(F3) are fixed proof concatenations and
numerical label tests.  Their fresh-label inequalities are quantifier-
free constructor inequalities, so their uniform correctness is
available in the weak base.

The inspected logical dependency boundary is exact.
In the 60 files of \file{PA196/Logic/}, primitive
\file{equalityVacuity} is consumed by
\file{Reverse.reverseChi} and
\file{Neighborhood.primitiveToNeighborhood}.  Replace the resulting
\file{reverseRight}, \file{reverseLeft}, and
\file{primitiveToNeighborhood} interfaces by (F2)--(F3).
\file{neighborhoodWitness} then uses (F3) and the unchanged witness.
The mirrored arithmetic certificate files have the same two consumers.
Other occurrences of \file{neighborhoodEqualityVacuity} in the
relational interface are derived from membership vacuity for the
defined relation $R$, as displayed in
\file{PA196/Logic/Relational/Interface.lean}; they are not primitive
equality-vacuity axioms.  The forward relational recipes above that
boundary consequently have exactly the required $\mathcal C_4$
implementations.  This is an interface replacement, not a purported
translation of arbitrary proofs using the obsolete axiom.

\subsection{Finite demand instead of induction over all environments}

For clarity the induction method is made explicit.  A request has a
rank $r$ and a finite parameter packet $x$.  Its finite child-query
list $Q(r,x)$ is computed before any child proof values are returned.
Same-stratum children have smaller rank.  A constructor $K$ combines
the supplied child lists, lower-stratum calls, and fixed leaves.
All validity and conclusion tests are primitive recursive.

For the binary structural recursions below, generate the complete
binary demand tree to depth $r+1$, retaining inactive placeholders.
There are at most $2^{r+2}-1$ slots.  A bounded forward loop computes
each active query from its parent; strict rank decrease proves that
all deepest active nodes are leaves.  Evaluate in reverse slot order.
This is primitive recursion, although individual parameter codes and
proof outputs can grow.  The output of the reverse loop itself gives
a primitive recursive bound on its intermediate outputs; no
preassigned small certificate bound is necessary.

The proof uses two finite Boolean invariants.  The first says that
all generated children of valid queries satisfy the stated local
validity and coverage tests.  The second says that every already
evaluated valid slot contains a checked list of its specified
conclusion.  Bounded conjunctions implement both as primitive
recursive Boolean equalities.  Quantifier-free induction proves them
with the root parameters free.  In $\mathsf I\Sigma _1$, the same
argument can equivalently use an existential finite witness table
with a bounded checking matrix.  It never induces on
$\forall\text{ environments}\,\exists\text{ proof}$.

Finite scans over an HFS set use an explicit supplied array of leaf
proofs, or a previously completed primitive recursive leaf function.
They do not turn an arbitrary $\forall i<k\,\exists e$ premise into
a primitive recursive selector.  Where a stratum calls a completed
lower stratum, the lower input rank may increase: this causes no
recursion cycle.  In particular vector-variable realization calls
term uniqueness on a replacement term which may be larger than the
original variable code; the companion term construction completes
uniqueness first.

\subsection{Coverage and the generic relational compiler}

Let $T_{b,n,\ell,z}(p)$ be the source relational translation, with
bound-label vector $b$, free-label vector $n$, allocator $\ell$, and
reserve $z$.  Its source constructors are translated to the displayed
numerical connectives; a quantifier has the child parameter map
\[
 A(b,n,\ell,z)=(\ell::b,n,\ell+1,z).
 \tag{F4}
\]
An equality atom is the defined relation $R$, with binder $z$;
a membership atom is $M$.  Function-free relational terms are bound
or free variables, so their atom labels are finite vector lookups.
This defines a primitive recursive translation by the finite-demand
construction.  It also supplies the same constructor equations as
the uniquely characterized upstream translation.

Write $\operatorname{Env}(b,n,\ell,z)$ for
$0<\ell\le z$ and every entry of $b,n$ being $<\ell$.
The exact coverage test is
\[
 \operatorname{Cov}(b,n,\ell,z,p)\ \equiv\
 \operatorname{Env}(b,n,\ell,z)\ \land\ \ell+p\le z
 \ \land\ \operatorname{IsSemiformula}_{|b|}(p)
 \ \land\ p\le|n|.
 \tag{F5}
\]
It implies coverage of conjunctive and disjunctive children at the
same parameters, and quantified children at (F4), because each child
code is strictly smaller than its parent.  The source lemmas are in
\file{FoundationRelationalCodeCoverage.lean} and already assume
only $\mathsf I\Sigma _1$.

For requested labels $\lambda=(n,\ell,z)$, abbreviate
$T_\lambda=T_{0,n,\ell,z}$ and put
\[
 P_d(\lambda)=(\operatorname{padEnvironment}(n,d),\ell,z+\ell+d).
 \tag{F6}
\]
Padding appends $d$ zero labels.  It preserves the entire old prefix,
has length $|n|+d$, and satisfies the same entry bound because
$0<\ell$.  For every source formula $q\le d$, (F6) supplies (F5)
whenever its syntactic arity is correct.  A finite vector loop proves
these facts with $n$ fixed.  Thus the full-PA induction over all
vectors in \file{FoundationRelationalSequentPadEnvironment.lean}
is replaced by this explicit loop.

The remaining generic relational operations have the following exact
same-family child queries.  Binary source connectives always have
their two children at unchanged ancillary parameters unless shown
otherwise; atoms and truth/false have no such child.
\begin{center}\small
\begin{tabular}{p{0.25\linewidth}p{0.64\linewidth}}
Operation & Quantifier child query for body $q$\\ \hline
NNF/native-negation comparison & $(q,A(b,n,\ell,z))$.\\
Low fresh-label vacuity & $(q,c,\ell::b,n,\ell+1,z)$, with $c<\ell$, $c\ne z$, and $c$ absent from $b,n$.\\
High-label vacuity & $(q,c,\ell::b,n,\ell+1,z)$, with $z<c$.\\
Conditional label renaming & $(q,x,y,\ell::b,\ell::\bar b,n,\bar n,\ell+1,z)$; corresponding vectors replace $x$ by $y$.\\
Reserve relocation & $(q,\ell::b,n,\ell+1,z,z')$.\\
Allocator relocation & With $m=\max(\ell,\ell')+1$, the two queries are $(q,\ell::b,n,\ell+1,m,z)$ and $(q,\ell'::b,n,\ell'+1,m,z)$.\\
Free-prefix agreement & $(q,\ell::b,n,n',\ell+1,z)$.\\
Free-variable shift agreement & $(q,A(b,n,\ell,z),x)$, where the other translation prepends $x$ to the free vector.\\
Bound substitution agreement & $(q,\operatorname{qVec}(w),\ell::a,\ell::b,n,\ell+1,z)$, starting from a label-substitution packet $(w,a,b,n)$.\\
\end{tabular}
\end{center}

For allocator relocation, the middle comparison is the already
completed conditional-renaming/binder-alpha constructor at allocator
$m$.  The two input coverage assumptions give $m+q\le z$; entries of
$\ell::b$ and $\ell'::b$ are $<m$.  The new eigenlabels are absent
from the old vectors because those entries were below both original
allocators.  No relocation call recurs through its output code.
For reserve relocation, the equality-atom case uses the finite
neighborhood binder-change recipe, and both reserves are distinct
from the atomic labels by (F5).  Membership atoms are unchanged.
Vacuity leaves use membership vacuity or the membership-derived
vacuity of $R$; the latter's reserve must also differ from the new
vacuity label, exactly as stated above.

The substitution packet condition is
$|w|=|a|$ and the label of $w_i$ in $(b,n)$ equals $a_i$ for every
$i<|w|$.  Under a binder it is preserved by
$(w,a,b)\mapsto(\operatorname{qVec}(w),\ell::a,\ell::b)$;
this is the exact \file{LabelSubstitution.qVec} calculation in
\file{FoundationRelationalSequentSubstitutionLabels.lean}.
Free-variable shift commutes with (F4), as recorded in
\file{FoundationRelationalSequentShiftLabels.lean}.
Prefix agreement assumes equality of entries at every $i<p$,
which covers all free indices occurring in $p$ and is inherited by
every smaller child.  Closed-environment independence can be proved
by a loop removing one free-vector entry at a time using the shift
equation and $\operatorname{shift}(p)=p$; it has fixed $p,b,\ell,z$
and a decreasing vector tail.  These are all Boolean syntax
equations.  The finite-demand invariant proves their uniform weak-
base versions even where an upstream theorem used an unbounded
all-environments induction.

The precise finite parent recipes are
\file{ActualCertificateEquivalenceInduction.source_step};
\file{FoundationRelationalSequentVacuity},
\file{HighVacuity}, \file{Rename}, \file{BinderAlpha},
\file{Allocator}, \file{Reserve}, \file{PrefixAgreement},
\file{Shift}, \file{Substitution}, and
\file{ClosedEnvironment}.  In this sentence the repeated prefix
\file{FoundationRelationalSequent} applies to each short name.
The proof uses their displayed case recipes with direct lists,
not their full-PA induction conclusions.

For an HFS sequent $s$, scan the numerical indices $i<s$ in ascending
order; membership is a bit test.  Start the disjunction fold at false.
At a member $i$, prepend the translated formula by disjunction;
at a nonmember, copy the previous result.  To eliminate the fold to
$C$, start with explosion $1\to C$ and at each member combine the
supplied arrow $T(i)\to C$ with the preceding arrow by the fixed
disjunction elimination recipe.  A member injection starts at its
own index using left injection and then uses right injection at each
later member.  Singleton, weakening, insertion, cut, and finite
premise discharge are these loops followed by fixed propositional
operations.  The arithmetic-formula sequent fold is identical with
leaf $T_\lambda(\operatorname{formulaSyntax}(e,i))$.
Target shift is a loop over the same HFS index set and the completed
formula-shift identity.  No identification of differently ordered
source-shifted scans is assumed; those are related by member
injection and fold elimination.

The generic relational LK compiler is now a finite-demand recursion
on its supplied derivation code $d$.  Its validity budget is
\[
 \operatorname{Env}(0,n,\ell,z),\qquad d\le|n|,
 \qquad \ell+d\le z.
 \tag{F7}
\]
The immediate proof children are strictly smaller than $d$.
Conjunction, disjunction, cut, and weakening pass $(n,\ell,z)$
unchanged; the universal rule passes $(\ell::n,\ell+1,z)$;
the existential rule passes $(n,\ell+1,z)$;
the shift rule passes $(\operatorname{tail}(n),\ell,z)$.
The last case is valid because $d>0$ and the child is $<d\le|n|$.
Identity and truth are fixed sequent constructors.
The axiom case calls the stated axiom supplier at the actual source
axiom and the current labels.  Formula coverage and the budgets for
the two raised-allocator rules follow from strict child inequalities
and (F7), exactly as in
\file{FoundationRelationalSequentDerivationStep.lean} and
\file{FoundationRelationalSequentProofBudget.lean}.

The universal parent recipe uses free-opening, allocator relocation,
and vacuity already completed above.  The existential parent recipe
uses covered term-label substitution, conditional renaming,
allocator relocation, and NNF comparison, as displayed in
\file{FoundationRelationalSequentExistentialInstance.lean} and
\file{FoundationRelationalSequentExsRule.lean}.
Thus this compiler depends only on lower relational strata and its
explicit axiom callback.  For the fixed-template uses below its
theory is empty, so that callback has no valid input.

\subsection{The nonrecursive leaves are supplied}

\paragraph{Fixed ND proof templates.}
An arbitrary formula hole in a fixed syntactic proof does not require
an arbitrary proof supplier.  In
\file{M6PA/Calibration/RelationalPredicateTemplateLKBounded.lean},
\file{template_lk_sound_of_bound} already assumes only
$\mathsf I\Sigma _1$.  The proof tree of each template is a fixed
finite host object.  Instantiate its formula holes, terms, shifts,
and binder extensions by primitive recursive code operations.  This
produces a source proof code $w$ without inspecting a numerical
certificate.  The fixed source-proof constructors are validated at
each host node; no induction over all internal formulas is used.

For a template with finite premises $\Delta$ and conclusion $q$,
let $s$ be its instantiated premise-set code and put
$b=\max(w,\max(s,q))$.  Compile $w$ by the generic relational
compiler at $P_b(\lambda)$, obtain each of the finitely many premise
proofs at that same label tuple, eliminate the negative-premise
sequent by the finite fold, and read $q$ back to $\lambda$ by
prefix agreement and reserve relocation.  The singleton form has
$b=\max(w,\max(a,q))$ for its premise $a$; the empty-premise form
uses the fixed truth proof.  These are exactly the budget formulas
in \file{RelationalPredicateTemplateNDPadded.lean},
\file{RelationalPredicateBinaryRequested.lean}, and
\file{RelationalPredicateBinaryEmptyRequested.lean}.
The child query labels depend on $w,a,q$, all computed from syntax,
and not on the returned premise-proof values.

This construction supplies the parameterized lower templates needed
by the companion term section: variable totality and coherence;
duplicate-domain and duplicate-both-domain; successor totality and
coherence; addition/multiplication totality and coherence;
substitution-variable coherence; the logical curry, pairing,
comparison, quantifier, and induction-combination templates.
Each is one of these fixed ND proof recipes, not a request to prove
an arbitrary formula.  Their only nonlogical premises are the
closed leaves in the next paragraph or the explicitly listed
callback supplies of the higher construction.

\paragraph{Closed relational leaves.}
The finite constant derivations needed here are the exact
\file{*NativeProof} definitions in the following files:
\begin{center}\small
\begin{tabular}{p{0.44\linewidth}p{0.45\linewidth}}
File & Supplied fixed conclusions\\ \hline
\file{RelationalClosedCertificateSeeds.lean} & inhabited zero, successor totality, addition totality, multiplication totality.\\
\file{RelationalClosedCertificateOrdinalSeeds.lean} & reflexivity, symmetry, congruence, transitivity, intersection, adjunction.\\
\file{RelationalClosedCertificateCoherenceSeeds.lean} & isomorphism reflexivity, symmetry, transitivity, addition coherence, multiplication coherence.\\
\file{RelationalClosedCertificateCoherenceZeroSeed.lean} & zero-isomorphism compatibility.\\
\file{RelationalClosedCertificateCoherenceTransportSeed.lean} & isomorphism transport.\\
\file{NumericPAeqSubstitutionVariableSeed.lean} & isomorphism right transport, from the explicit fixed \file{isoRightTransportDerivation}.\\
\file{RelationalClosedCertificateCoherenceSuccessorSeed.lean} & successor-isomorphism compatibility.\\
\file{PAInterpretedFixedCertificateSeeds.lean} & the twelve fixed ordinary-PAeq axiom images.\\
\end{tabular}
\end{center}
These are finite source proof recipes at fixed standard reserves,
with numerical reserves explicitly larger than every listed code and
LK budget.  For example
\begin{align*}
 \operatorname{primitiveReserve}(q)
  ={}&\operatorname{functionBudgetBound}(1)\\
     &{}+(1+\operatorname{lkBudget}(\mathrm{mulUniversalOutputDerivation}))\\
     &{}+(1+\operatorname{encode}(q)),
\end{align*}
and \file{coherenceReserve} is the displayed maximum of its
thirteen local budgets plus $1+\operatorname{encode}(q)$.
The remaining fixed reserves use
\file{certificateReserve} with the displayed two inputs.  There
is consequently no search for a sufficient reserve.

Reinterpret these finite derivation recipes through the logical
interface (F2)--(F3).  The native relational source compiler occurring
inside them traverses a fixed host proof tree; its finite
separation calls are instances of the retained gate construction.
All numerical side conditions are the explicit budget inequalities
in the definitions.  This gives a fixed $\mathcal C_4$ proof list
for each closed conclusion.  Its validity is a finite arithmetic
calculation.  We have inspected the definitions and their interface
boundary; we do not claim to have evaluated their very large
numeral codes or replayed a new Lean build.

To supply such a closed conclusion at arbitrary covered requested
labels, start at its recorded zero-label environment, allocator $1$,
and fixed reserve; use closed-environment independence, allocator
relocation, and reserve relocation proved above.  This is the direct
list interpretation of
\file{RelationalClosedCertificate.relocated_certificate}.
For the twelve ordinary fixed axioms, use the finite case split of
\file{PAInterpretedFixedNumericClass.lean}; its conclusion is the
literal two-stage \file{formulaSyntax} image, whose identity with
the closed fixed source formula is the displayed constructor
calculation in the seed-calibration/relocation files.  No arithmetic
soundness or truth of PA is being assumed by these fixed leaves.

\paragraph{The arbitrary-property separation leaf.}
This is the only nonconstant axiom supplier required by the ordinal
induction template.  For a valid relational one-variable formula
$p$, the requested source separation formula has property
translation
\[
 B=T_{(\ell+2)::0,n,\ell+3,z}(p).
\]
Let $a=\ell$, $d=\ell+2$, and $c=\ell+1$.  Coverage of the source
separation formula gives the stated property coverage, makes $c$
fresh for its vectors, and puts all three labels below $z$.
Put
\[
 q=U(z+1,M(a,d)\wedge B)>z.
 \tag{F8}
\]
High-label vacuity of $B$ at $z+1$, and membership vacuity of
$M(a,d)$ there, give the padding equivalence used by
\file{paddedSeparationAbove}.  Pointed separation and the old gate
give the raw separation conclusion with output name $q$.
The already supplied high vacuity at $q$ and low fresh vacuity at
$c$ feed \file{separationOutputAlpha}; (F3), membership rewrites,
and its explicit inequalities change the output name to $c$.
Gen at $\ell$ and the fixed three-quantifier normalization read this
as the exact source separation schema.  The source equations and
recipes are in
\file{FoundationRelationalSequentSeparationSupplier.lean},
\file{SeparationAtName.lean}, \file{SeparationReadback.lean},
\file{RawSeparation.lean}, and
\file{ActualCertificateSeparationRaw.lean}/
\file{ActualCertificateSeparationRelocation.lean}
(with the full prefix on the four shortened sequent filenames).
Only the already completed structural vacuity and syntax algorithms
recurse on $p$.  No PA formula-realization or PA-proof compiler is
called here.

\subsection{Requested-label composition and the arithmetic syntax}

To avoid confusing source codes with the numerical implication $I$,
write $[A\Rightarrow B]$ for the source code
\file{imp language A B}, and $A\mathbin\&B$ for source conjunction.
The exact padding budgets are
\begin{align*}
 \kappa(A,B,C)&=\max([A\Rightarrow B],
                         \max([B\Rightarrow C],[A\Rightarrow C])),\\
 \eta(A,B,C)&=\max([A\Rightarrow B],
                         \max([A\Rightarrow C],[A\Rightarrow B\mathbin\&C])),\\
 \mu(A,B)&=\max(A,\max([A\Rightarrow B],B)).
 \tag{F9}
\end{align*}
Requested composition and pairing query their children at
$P_\kappa(\lambda)$ and $P_\eta(\lambda)$ respectively.
Requested MP queries both inputs at $P_\mu(\lambda)$.
Native logical composition is followed by source/native connective
normalization and padding readback.  All three intermediate source
formula codes are bounded by the displayed maximum, so (F5)--(F6)
prove every child coverage requirement.

For a one-variable source formula $p$, requested universal closure
asks for its opened formula at
\[
 O_p(\lambda)=
 (\ell::\operatorname{padEnvironment}(n,d),\ell+1,z+\ell+d),
 \qquad d=\operatorname{Free}(p)+1.
 \tag{F10}
\]
The source equations for \file{Free} proved above identify the
body at allocator $\ell+1$ exactly.  It is not identified at the
parent allocator.  Gen, allocator relocation, and readback finish
the requested universal formula.  The budget and query are exactly
\file{FoundationRelationalSequentRequestedUniversal.lean}.

Two derived opening operations used below are specified by these
same maps.  Guarded universal closure of $[H\Rightarrow\forall P]$
is MP with antecedent $\forall[H\Rightarrow P]$, whose first input
is requested universal closure of $[H\Rightarrow P]$ and whose
second is the fixed universal-curry template.  Its child is therefore
queried at
$O_{[H\Rightarrow P]}(P_{\mu(\forall[H\Rightarrow P],
 [H\Rightarrow\forall P])}(\lambda))$.
Double universal closure of $\forall\forall P$ is two requested
universal operations, with child labels
$O_{\operatorname{openOuter}(P)}(O_{\forall P}(\lambda))$,
where
$\operatorname{openOuter}(P)=
 \operatorname{subst}(\operatorname{qVec}(\operatorname{fvar}(0)::0),
                      \operatorname{shift}(P))$.
The exact identities are
$\operatorname{Free}(\forall P)=\forall\operatorname{openOuter}(P)$
and its once-more opened body, as displayed in
\file{NumericPAeqFormulaDoubleUniversal.lean}.

The arithmetic value-syntax construction and its substitution,
shift, validity, and support joins are supplied by the companion
term section.  Arithmetic \file{formulaSyntax} is then another
finite-demand syntax construction: constants are unchanged;
positive/negative equality atoms use the two value-syntax outputs;
binary connectives recurse on the two children with unchanged
parameters; an arithmetic quantifier recurses on its body at
\file{binderParameters} and wraps the result with the fixed
domain-guarded universal/existential constructor.  These are the
literal cells of
\file{M6PA/Calibration/NumericPAeqFormulaConstruction.lean}.
The term isomorphism/domain formulas in these fixed constructors
are relational syntax, not numerical proof oracles.

The formula support and substitution joins are proved by the same
finite-demand tree: binary cases use unchanged parameters,
quantifiers use \file{binderParameters} and lifted substitution
$\operatorname{qVec}$, atoms use the completed term substitution
join.  In the target-substitution join, a binder increments both
target arities and lifts the target vector by $\operatorname{qVec}$;
its two parameter packets agree by the explicit shift/substitution
commutation.  Closed free-environment independence follows because
the support bound is zero.  Thus the following uses of
\file{formulaSyntax} have actual primitive recursive meanings and
weak-base constructor equations; they do not inherit a
\file{Classical.choose} totality premise from a stronger model.

\subsection{Arithmetic formula realization: the actual child requests}

Let $e=(s,r,b,f,a,w)$ satisfy the exact
\file{FormulaEnvironment} conditions: $b$ has $r$ valid relational
terms of arity zero; $a$ has $s$; $f$ has $|f|$; the arithmetic source
formula $p$ has arity $s$; $w$ is a source term vector of length $s$
and target arity $r$; $p$ and every entry of $w$ have free support
below $|f|$.  These are all bounded syntactic tests.  Define
\[
 H_e=\operatorname{vectorInput}(s,r,b,f,a,w),
\]
\[
 C_e(p)=\operatorname{Bi}\bigl(
 \operatorname{formulaSyntax}(b,f,\operatorname{subst}(w,p)),
 \operatorname{formulaSyntax}(a,f,p)\bigr),\qquad
 L_e(p)=[H_e\Rightarrow C_e(p)].
\]
Here $\operatorname{Bi}(A,B)=[A\Rightarrow B]\mathbin\&[B\Rightarrow A]$.
The requested realization builder returns a proof of
$T_\lambda(L_e(p))$ under its coverage test.

For an atomic equality of terms $t,u$, put $H=H_e$ and let the six
conclusions be the two vector comparisons for $t,u$, followed by
the four totality formulas for the substituted $t,u$ and original
$t,u$, in that order.  Call them $B_1,\ldots,B_6$ and set
$R_i=B_i\mathbin\&R_{i+1}$, with $R_6=B_6$.
The atom-input supplier is the right-nested
$\operatorname{Pair}(H,B_1,R_2;\ldots)$ through
$\operatorname{Pair}(H,B_5,B_6;\ldots)$.
Compose it with the fixed equality-transport template using
$\operatorname{Comp}(H,R_1,C_e(t=u);\ldots)$.
This specifies all six query labels by (F9).
The first two leaves invoke completed vector realization.  Each
substituted totality leaf is
$\operatorname{Comp}(H,\operatorname{inputDomain}(r,b,f),B_i;
 \operatorname{LeftProjection},\operatorname{Tot}(\operatorname{subst}(w,t)))$
(or $u$); each original totality leaf uses
$\operatorname{inputDomain}(s,a,f)$ and original $t$ or $u$.
That domain projection is itself the two fixed conjunction
projections displayed in
\file{NumericPAeqSupportedVectorRealizationTotality.lean}.
Thus all six leaves belong to completed term strata even if a
substituted term is large.
Negative equality is composition of this positive-atom helper with
the fixed negated-comparison template, at its $P_\kappa$ labels;
it makes no recursive formula-supply query.  Truth and false are
fixed logical templates.  The exact recipes are
\file{NumericPAeqFormulaRealizationAtomInput.lean},
\file{NumericPAeqFormulaRealizationAtomTransport.lean}, and
\file{NumericPAeqFormulaRealizationNegativeAtom.lean}.

For a binary parent $p\circ q$, with $\circ$ either source
conjunction or disjunction, set
$A=C_e(p)$, $B=C_e(q)$, $C=C_e(p\circ q)$.
The two recursive requests are exactly
\[
 (p,e,\lambda_2),\quad(q,e,\lambda_2),\qquad
 \lambda_1=P_{\kappa(H_e,A\mathbin\&B,C)}(\lambda),\quad
 \lambda_2=P_{\eta(H_e,A,B)}(\lambda_1).
 \tag{F11}
\]
Pair the two arrows and compose with the fixed conjunction/disjunction
comparison template.  These are the case recipes of
\file{NumericPAeqFormulaRealizationConnectives.lean}.

For a universal or existential parent with body $p$, define the
source one-variable matrix
\[
 W=[H_e\Rightarrow[D_0\Rightarrow
            \operatorname{Bi}(\operatorname{binderLeft}(b,f,w,p),
                              \operatorname{binderRight}(a,f,p))]],
\]
where $D_0$ is the fixed domain template.  Let
$J=L_e(\forall p)$ or $L_e(\exists p)$ as appropriate, and
$x=\operatorname{fvar}(0)$.  Put
\[
 e^+=(s+1,r+1,x::\operatorname{sh}(b),\operatorname{sh}(f),
                   x::\operatorname{sh}(a),\operatorname{qVec}(w)),
\]
where $\operatorname{sh}$ shifts each relational target term.
Also put
\[
 H=\operatorname{vectorInput}(s,r,\operatorname{sh}(b),
       \operatorname{sh}(f),\operatorname{sh}(a),w),\quad
 D=\operatorname{domainCode}(x),\quad I_0=H_{e^+},\quad C=C_{e^+}(p).
\]
The sole recursive formula request is $(p,e^+,\lambda_4)$, with
\begin{align*}
 \lambda_1&=P_{\mu(\forall W,J)}(\lambda),&
 \lambda_2&=O_W(\lambda_1),\\
 \lambda_3&=P_{\mu([H\mathbin\&D\Rightarrow C],
                           [H\Rightarrow[D\Rightarrow C]])}(\lambda_2),&
 \lambda_4&=P_{\kappa(H\mathbin\&D,I_0,C)}(\lambda_3).
 \tag{F12}
\end{align*}
At the last line use the companion term builder's explicit
\file{vector_qVec_input_requested_certificate} for
$H\mathbin\&D\to I_0$, then the recursive formula child,
then the fixed curry, universal-closure, and universal/existential
comparison templates.  The literal opening identity is
\[
 \operatorname{Free}(W)=[H\Rightarrow[D\Rightarrow C]].
\]
It is exactly \file{uniformMatrix_free} in the binder-opening
file, obtained from the completed source and target substitution
joins.  The environment conditions for $e^+$ follow by lifting
$w$, shifting the three target vectors, and prepending $x$ to $a,b$.
They are the explicit bounded checks of
\file{environment_targetShift} and
\file{environment_all_body}; (F9)--(F10) supply all coverage
checks.  Since $p$ is smaller than either quantified parent,
(F11)--(F12) and the supplied leaves close the finite-demand
construction for \emph{every} arithmetic formula.

\subsection{The arbitrary induction-matrix supplier}

For a relational predicate code $P$ of arity one, the fixed ordinal
induction proof template has exactly seven premises.  Their precise
dispatch is \file{RelationalPredicateTemplateOrdinalCertificate.}
\file{induction_axiom_certificate}: reflexivity, symmetry,
congruence, transitivity, intersection, adjunction, and source
separation for $\neg P$.  The first six are the fixed closed leaves
above; the seventh is (F8).  The template compiler has already been
constructed.  Thus the ordinal-induction supplier is total with
primitive recursive witnesses for arbitrary $P$, and it has no
call back to arithmetic formula realization or PA-proof transport.

For an arithmetic induction property $p$ and target free-term vector
$f$, the canonical induction-body constructor is the exact finite
composition in \file{NumericPAeqInductionBodyCertificate.lean}.
Using its displayed formulas $F,B,S,A,T,C$, define
\[
 X=[F\mathbin\&B\Rightarrow A],\quad
 Y=[F\mathbin\&S\Rightarrow T],\quad
 Z=[A\mathbin\&T\Rightarrow C],\quad
 R=[F\Rightarrow[B\Rightarrow[S\Rightarrow C]]].
\]
Its request tree, naming the two inner subexpressions, is
\begin{align*}
 C_X&=\operatorname{MP}(X,[Y\Rightarrow[Z\Rightarrow R]];
          \operatorname{BaseCompatibility},\operatorname{Combine}),\\
 C_Y&=\operatorname{MP}(Y,[Z\Rightarrow R];
          \operatorname{StepCompatibility},C_X),\\
 C_Z&=\operatorname{MP}(Z,R;\operatorname{OrdinalInduction},C_Y).
 \tag{F13}
\end{align*}
Each MP includes exactly its $P_\mu$ query from (F9).
Base compatibility is a requested guarded universal closure,
followed by MP of the zero-instance formula comparison and the fixed
instance-forward template.  Step compatibility is MP of a double
requested universal closure of the successor comparison and the
fixed step template.  Zero-instance and successor comparison invoke
the already completed formula-realization builder and term
suppliers.  The finite recipe files are
\file{NumericPAeqInductionBaseCompatibility.lean},
\file{NumericPAeqInductionStepCompatibility.lean},
\file{NumericPAeqInductionSuccessorComparison.lean},
\file{NumericPAeqFormulaZeroInstance.lean}, and
\file{NumericPAeqFormulaSuccessorInstance.lean}.
Their parameters are source formula/term codes and the displayed
shifted vectors; none is a returned numerical proof value.

For a canonical matrix $b$ with $m$ bound variables and witness
$\operatorname{subst}(\operatorname{fvarVec}(m),b)
 =\operatorname{targetIndBody}(p)$, let $B$ be an arbitrary valid
target-term vector of length $m$.  Use the notation of
\file{NumericPAeqCanonicalMatrixCertificate.lean}:
\[
 G=\operatorname{guard}(B,m),\quad
 I_0=\operatorname{vectorInput}(m,0,0,B,B,\operatorname{fvarVec}(m)),
\]
\[
 A=\operatorname{substitutedFormula}(0,B,\operatorname{fvarVec}(m),b),
 \qquad R=\operatorname{originalFormula}(B,B,b).
\]
The actual constructor, naming its inner subexpressions, is
\begin{align*}
 C_0&=\operatorname{Comp}(G,I_0,\operatorname{Bi}(A,R);
                 \operatorname{CanonicalInput},\operatorname{FormulaRealization}),\\
 C_1&=\operatorname{Pair}(G,A,\operatorname{Bi}(A,R);
                 \operatorname{InductionBody},C_0),\\
 C_2&=\operatorname{Comp}(G,A\mathbin\&\operatorname{Bi}(A,R),R;
                 C_1,\operatorname{UseBi}).
 \tag{F14}
\end{align*}
All queries in (F14) include their (F9) label maps.  The canonical
input proof uses the completed explicit vector/guard construction.
The source identity of the induction-body antecedent follows from
the supplied substitution equation; closed free-support identifies
$R$ with the desired matrix.  Projecting
$\operatorname{inputDomain}(m,B,f)$ to $G$ supplies arbitrary $f$.
This is a completed supplier for $b$, not an induction hypothesis
asserting supplies for all environments.

For universal closure first fix a canonical witness packet
$(m,p,K)$ with
\begin{gather*}
 \operatorname{IsSemiformula}_m(p),\quad
 \operatorname{shift}(p)=p,\quad \operatorname{bv}(p)=m,\quad
 \operatorname{IsSemiformula}_1(K),\\
 \operatorname{subst}(\operatorname{fvarVec}(m),p)
        =\operatorname{targetIndBody}(K).
\end{gather*}
Every request carries this entire packet unchanged.  Define a
request $\mathcal M(m,p,K;k,s,b,f,\lambda)$, whose validity test
also requires $s+k=m$, and whose output is a proof of
\[
 T_\lambda\bigl(\operatorname{matrixCode}
                   (s,b,f,\operatorname{Alls}(p,k))\bigr).
\]
The packet gives source arity $s+k$ and no free variables.  At
$k=0$ the invariant gives $s=m$, and (F14) applies with the
recorded witness $K$.  This base case does not supply arbitrary
closed matrices.
For the $k+1$ case put $q=\operatorname{Alls}(p,k)$,
$H_0=\operatorname{inputDomain}(s,b,f)$,
\[
 P_0=[D_0\Rightarrow
       \operatorname{formulaSyntax}(\operatorname{binderParameters}(b,f),q)],
 \quad A_0=\forall[H_0\Rightarrow P_0],\quad
 B_0=[H_0\Rightarrow\forall P_0].
\]
Let $x=\operatorname{fvar}(0)$, $b'=x::\operatorname{sh}(b)$,
$f'=\operatorname{sh}(f)$, and
\begin{align*}
 H&=\operatorname{inputDomain}(s,\operatorname{sh}(b),f'),&
 D&=\operatorname{domainCode}(x),\\
 I_0&=\operatorname{inputDomain}(s+1,b',f'),&
 Q&=\operatorname{formulaSyntax}(b',f',q).
\end{align*}
The only same-family child is
$\mathcal M(m,p,K;k,s+1,b',f',\lambda_4)$, where
\begin{align*}
 \lambda_1&=P_{\mu(A_0,B_0)}(\lambda),&
 \lambda_2&=O_{[H_0\Rightarrow P_0]}(\lambda_1),\\
 \lambda_3&=P_{\mu([H\mathbin\&D\Rightarrow Q],
                         [H\Rightarrow[D\Rightarrow Q]])}(\lambda_2),&
 \lambda_4&=P_{\kappa(H\mathbin\&D,I_0,Q)}(\lambda_3).
 \tag{F15}
\end{align*}
The unchanged witness packet and $s+1+k=m$ validate this child.
The other inputs are the fixed input-prepend, curry, and universal-
curry templates.  These are the displayed cases of
\file{NumericPAeqGuardedUniversalCertificate.lean} and
\file{NumericPAeqGuardedAlls.lean}.  Rank $k$ strictly decreases
and every parameter is calculated before the child proof is returned.
The finite-demand proof therefore supplies the exact replacement
for \file{matrixSupply_alls}.

Finally, \file{targetInductionR}(q) has bounded witnesses
$m,b\le q$ and an induction property
$p\le\operatorname{subst}(\operatorname{fvarVec}(m),b)$.
The remaining tests are source syntax, $q=\operatorname{Alls}(b,m)$,
$\operatorname{shift}(b)=b$, $\operatorname{bv}(b)=m$, and the
displayed opened-body equality.  Decode it by bounded search over
these bounds.  The last bound is on the computed opened-body code;
it is not silently replaced by $q$.  Thus the induction-axiom
dispatcher has a primitive recursive witness extractor.
Apply (F15) at $k=m,s=0$ and then MP with the fixed proof of
$\operatorname{inputDomain}(0,0,0)=\top\mathbin\&\top$.
Closed free-support removes the arbitrary vector $f$.  This supplies
the literal \file{formulaSyntax} image of every ordinary induction
axiom at every covered requested label tuple.

\subsection{The outer ordinary-PAeq proof compiler}

Write $S_f(s)$ for the arithmetic-formula sequent fold at source
term environment $f$ and
$J_f(s)=[\operatorname{freeDomain}(f)\Rightarrow S_f(s)]$.
A request $(d,f,\lambda)$ assumes that $d$ is a valid ordinary-PAeq
derivation, $f$ is a vector of valid arity-zero relational terms,
$d\le|f|$, and $J_f(\operatorname{fstIdx}(d))$ is covered at
$\lambda$.  The output is a direct proof of that translated
guarded sequent.  The proof-child code is the recursion rank.

The exact requested-child maps are as follows.  Here $t,u$ denote
the actual child sequents; they are read from the proof node, not
chosen by proof search.
\begin{center}\small
\begin{tabular}{p{0.18\linewidth}p{0.71\linewidth}}
Rule & Recursive request(s)\\ \hline
Identity or truth & None; NNF normalization and the finite sequent member templates.\\
Conjunction or cut & The two actual proof children, with the same $f$, at $P_{\max(J_f(t),J_f(u))}(\lambda)$.\\
Disjunction or weakening & Its single proof child at $P_{J_f(t)}(\lambda)$; the binary-lift recipe can reuse that one output.\\
Existential & Its actual proof child with $f$ unchanged at $P_{\max(J_f(t),\operatorname{existentialRuleCode}(f,p,v))}(\lambda)$.\\
Shift & Write $f=x::f_0$; request the proof child with $f_0$ at $P_{J_{f_0}(t)}(\lambda)$.\\
Axiom & No proof child; request the fixed/induction axiom supplier at $P_{\operatorname{formulaSyntax}(0,f,p)}(\lambda)$.\\
Universal & The single proof child with $f'=x::\operatorname{sh}(f)$ at $\lambda_6$ defined below.\\
\end{tabular}
\end{center}
These are the literal rules in
\file{NumericPAeqSequentDerivationStep.lean},
\file{RequestedLift.lean}, \file{RequestedLeaves.lean},
\file{ExsRule.lean}, and \file{ShiftRule.lean}, with the prefix
\file{NumericPAeqSequent} on each shortened filename.

For the universal rule let $p$ be its body, $s$ its parent sequent,
$H_0=\operatorname{freeDomain}(f)$, $S=S_f(s)$,
$A=\operatorname{formulaSyntax}(0,f,\forall p)$,
$J=A\lor S$, and
\[
 W=[H_0\Rightarrow[D_0\Rightarrow
          (\operatorname{property}(f,p)\lor S)]],\qquad
 R=[H_0\Rightarrow J].
\]
Put $x=\operatorname{fvar}(0)$, $f'=x::\operatorname{sh}(f)$,
$H=\operatorname{freeDomain}(\operatorname{sh}(f))$,
$D=\operatorname{domainCode}(x)$,
$I_0=\operatorname{freeDomain}(f')$, and
\[
 C=\operatorname{Free}(\operatorname{property}(f,p))
                   \lor\operatorname{shift}(S).
\]
Then the actual child labels are
\begin{align*}
 \lambda_1&=P_{\kappa(H_0,J,S)}(\lambda),\\
 \lambda_2&=P_{\mu(\forall W,R)}(\lambda_1),\\
 \lambda_3&=O_W(\lambda_2),\\
 \lambda_4&=P_{\mu([H\mathbin\&D\Rightarrow C],
                       [H\Rightarrow[D\Rightarrow C]])}(\lambda_3),\\
 \lambda_5&=P_{\kappa(H\mathbin\&D,I_0,C)}(\lambda_4),\\
 \lambda_6&=P_{J_{f'}(\{\operatorname{Free}(p)\}\cup
                          \operatorname{setShift}(s))}(\lambda_5).
 \tag{F16}
\end{align*}
The finite source-sequent shift/fold templates identify the actual
child context with this formula.  The universal-context recipe then
uses input-prepend, curry, and its fixed uniform template.  This is
\file{NumericPAeqUniversalContextCertificate.lean},
\file{NumericPAeqUniversalContextOpening.lean}, and
\file{NumericPAeqSequentAllRule.lean}.

All proof children are $<d$.  The source proof-budget conditions
descend because the unchanged vector has length at least $d$;
the eigenvector gains one entry; the shift tail loses one entry
but its child is strictly smaller.  A source term in an existential
node has code below its parent proof code, so its free support is
covered.  Every requested native formula has its needed code bound
from (F9), (F10), or the explicit maxima in the table.  The target
syntax/support identities are the completed finite syntax
constructions.  These facts prove the local coverage invariant
without an all-environments induction.

The existential parent uses an arbitrary-term totality proof and
the completed formula realization for its source substitution,
then its fixed existential template.  The axiom parent invokes the
finite fixed-axiom case split or the completed (F13)--(F15) induction
supplier.  Neither calls this ordinary-proof compiler.  Thus the
finite-demand algorithm on $d$, using the exact table and (F16),
defines a total primitive recursive requested-sequent translator.
The reverse evaluation invariant proves its correctness in the
weak base.

\subsection{Discharging the domain guard and proving (F1)}

Given an ordinary-PAeq refutation code $d$, choose the relational
source environment $f$ of length $d$ all of whose entries are
$x=\operatorname{fvar}(0)$.  It satisfies the source proof budget.
The outer compiler supplies the guarded singleton-false sequent.
The explicit repeated-domain loop in the companion term section
gives
$\operatorname{domainCode}(x)\to\operatorname{freeDomain}(f)$:
its zero case is truth weakening and its successor case pairs the
previous arrow with identity.  Thus it does not need an unspecified
array of domain witnesses.

The singleton sequent fold and requested composition give
$\operatorname{domainCode}(x)\to\bot$.
Requested universal closure gives
$\forall x(\operatorname{domainCode}(x)\to\bot)$.
The fixed closed inhabited-zero leaf is
$\exists x(\operatorname{domainCode}(x)\land\operatorname{zeroCode}(x))$.
The finite \file{domainContradiction} template and two requested
MP operations yield the source false formula, whose numerical
translation is literally $1$.
These are the exact joins in
\file{NumericPAeqBottomGuardDischarge.lean},
\file{NumericPAeqBottomDomainCertificate.lean}, and
\file{NumericPAeqRefutationTransport.lean}.
Start the requested final false formula at a zero-label vector of
length its source code, allocator $1$, and reserve one plus that
code; all (F5) conditions hold.  Expanding the displayed requested
composition/MP/universal operations fixes every upstream request
label by (F6), (F9), and (F10).  No choice is left to a proof oracle.

Define $F(d)$ to be this composition of the finite-demand builders,
with zero as its arbitrary default output when an input validity
test fails.  Its complete stratum order is:
\begin{enumerate}
\item primitive source syntax, list checking, and the finite logical
      recipes including (F2)--(F3);
\item source translation, syntax joins, NNF, vacuity, renaming,
      binder alpha, allocator/reserve relocation, and finite HFS folds;
\item generic relational LK, finite template instantiation/discharge,
      the fixed closed seeds, and arbitrary-property separation;
\item arithmetic value/formula syntax and its finite substitution
      joins; term projections, totality, coherence, uniqueness, and
      vector realization as in the companion term section;
\item arithmetic formula realization (F11)--(F12);
\item ordinal induction, canonical induction body and matrix
      (F13)--(F14), universal matrix closure (F15), and ordinary
      axiom dispatch;
\item the ordinary proof compiler and its universal child map (F16),
      followed by the explicit bottom/domain discharge.
\end{enumerate}
The few syntax constructions used in an earlier item can be moved
before all proof strata: their recursive calls are entirely on
syntax and have no certificate supplier.  Within every genuine
proof recursion the rank decreases along the displayed same-family
edges; every other edge goes to a previously completed stratum.
All variable-length leaves are bounded scans or term/formula
recursions already listed.  The remaining leaves are fixed recipes,
with all their assumptions supplied above.  This proves (F1).

\paragraph{What is and is not imported from the stronger development.}
The original arbitrary-formula realization theorem and
\file{matrixSupply_alls} used full-PA induction on unbounded
families of environments; other source relocation and sequent
theorems did likewise.  The proof here imports their explicit
constructor equations and finite parent proof recipes, replaces
their noncomputable witness interfaces by the displayed algorithms,
and proves correctness on finite query tables.  It does not cite
those stronger induction conclusions as weak-base lemmas.  The
fixed standard native proofs are cited as finite derivation data
after the exact diagonal-interface replacement.  That distinction
is what permits a weak-base calibration while leaving the upstream
Lean status unchanged.

For $\mathcal C_{4Q}$ replace $h_6$ by its antecedent-exchanged head.
The finite restricted-propositional exchange constructor above
recovers either order from the other.  Replacing each such head
row by that finite proof-list expansion gives the corresponding
primitive recursive translator and the same theorem (F1) with
$\operatorname{C4QProof}$.
